\section*{Introduction}

Potentially harmful misinformation about social, political, and public health issues is a growing problem. The problem is particularly acute in the Global South, where digital media literacy remains low and social media is often the only affordable source of online information \citep{arechar2022understanding}. The COVID-19 infodemic, which has likely led to harmful individual health choices around the world, highlights the need to identify scalable ways to counter the spread of misinformation and its effects. Indeed, conspiracy theories surrounding the pandemic and the vaccines may play an important role in contributing to low vaccination rates across the Global South %remain low partly because of conspiracy theories surrounding the pandemic and the vaccines 
\citep{mallapaty2022researchers, ackah2022covid}.

Existing research highlights the utility of both \textit{debunking} and \textit{prebunking} interventions \citep{Cooketal2017Inoculation, Nyhanetal2020Literally, walter_meta, guess2020digital, Vragaetal2020}, but each mode of combating misinformation faces major obstacles to implementation. Debunking, which principally involves providing fact-checks designed to counteract misinformation \citep{vanderMeerYan2020Formula, nyhan_reifler_2015}, faces two connected problems. First, fact-checks are usually released too late: fact-checking occurs when the virality of the misinformation on social media is already naturally decaying and misinformation has already caused harm. Second, fact checks have very limited virality relative to the misinformation they aim at debunking, so they fail to reach the wide audience affected by misinformation. 

Prebunking, which instead relies on literacy training and warnings as a way of inoculating people against misinformation, seeks to address the timing and virality issues by attempting to prevent belief in misinformation \textit{before} citizens encounter misinformation \citep{Compton2013Inoculation, mcguire1964some}. However, while prebunking interventions have demonstrated effectiveness in various settings, they still involve recruiting and incentivizing people to consume information. It is unclear whether participants will continue to engage once they are no longer incentivized by an intervention, whether positive results are long-lasting, and whether these interventions can be scaled at low cost.

We tackle these outstanding issues. The study  was developed in partnership with our institutional collaborator, Africa Check, the first and largest fact-checking organization in sub-Saharan Africa, and it builds on previous work with Africa Check, Chequea Bolivia, one of the two largest fact-checkers in Bolivia, and the independent media development organization Internews. Our goal is to inform the programming decisions that fact-checkers face relating to mitigating misinformation's potentially deleterious effects. 

Specifically, we ask whether social media micro-influencers can be recruited to increase awareness and knowledge of misinformation among their online audiences. The effectiveness of fact-checking interventions relies heavily on users' trust in and attention to a source \citep{bowlesetal2020}. Consequently, social influencers with large and engaged audiences may be an especially effective mode of delivery to target large numbers of individuals at a minimal cost. However, this raises a critical question: can credible social influencers be incentivized to disseminate information about misinformation? While piloting the recruitment process with Chequea Bolivia indicates that financial incentives for posting fact-checks are required for compliance (in the form of compensating the fact-checkers for internet data costs and time), their effectiveness is unclear. 

Moreover, we compare the effectiveness of recruiting social media micro-influencers to alternatively using the funds to target ads to social media followers. One advertising is logistically simpler and has the appeal of generating significantly more impressions for the same cost. However, its lower trustworthiness might mean it has a lower ultimate impact. 

The experiment has been implemented with 302 influencers across Kenya and South Africa in three phases: \textit{pilot}, \textit{first batch}, and \textit{second batch}. The \textit{pilot} was implemented between December 5th, 2022 and January 29th 2023, the \textit{first batch} between March 13th, 2023 and May 7th, 2023, and the \textit{second batch} started on May 1st, 2023 and ended on June 25th, 2023.

We focus on three main sets of outcomes only for the first and second batch phases, which we measure for the followers of social media influencers. First, we are interested in whether social media usage patterns change as a result of our intervention. Specifically, we measure the extent to which the followers of treated social media micro-influencers become more likely to post and engage with (1) content of higher quality, (2) content that reports fact-checks, and (3) content that misinforms the readers. To that end, we have scraped Twitter posts and used machine/deep-learning techniques we have started to develop in the context of \cite{Bandieraetal2023} to classify whether the posts are verifiable, and, if so, whether they constitute likely misinformation (see more details in Appendix \ref{ver_bert}). We also investigate using Natural Language Processing Techniques (NLP) how treatment changes engagement with social media content about other topics usually affected by misinformation, including health, politics, and economics. 

Second, through surveys, we assess whether the first intervention affects trust and consumption of different media sources, the ability to discern between true and false news, news sharing and verification behavior, as well knowledge, sentiments, and behaviors with regard to vaccine uptake. 

Regarding data collection, first, we scraped baseline Twitter posts for the pilot and first and second batch influencers' followers (3 months before the start of each intervention). Additionally, we scraped the Twitter posts for all eight weeks of the first-batch followers and some of the second-batch followers. Lastly, we scraped one month of post-intervention Twitter posts for the first and some followers of the second batch. We are in the process of scraping Twitter posts for the intervention and post-intervention periods for the remainder of the second-batch followers.

%Regarding data collection, first, we have scraped  baseline Twitter posts for the first and second batch influencers' followers (3 months before the start of each intervention). Additionally, we have scraped the Twitter posts for all eight weeks of the first-batch and second-batch interventions. Lastly, we have scraped one month of post intervention Twitter posts for the first and second batches. \footnote{We are  missing the baseline Twitter posts for the second batch, some of the  \textit{second-batch} intervention period tweets, and the post-intervention period for the \textit{second-batch}} Second, we have collected survey data for the \textit{first batch} intervention and the \textit{second batch} intervention. To obtain this data, we utilized the Academic Twitter API to send direct messages to all influencers' followers who had their messages open, inviting them to participate in a brief survey in exchange for a chance to win a monetary prize.
